% ----------------------------------------------------
% Conclusion
% ----------------------------------------------------
\documentclass[class=report,11pt,crop=false]{standalone}
\input{../Style/ChapterStyle.tex}
\begin{document}
% ----------------------------------------------------
\chapter{Conclusion} \label{ch:conclusion}
\vspace{-1cm}
% ----------------------------------------------------
% This CHAPTER only needs to be completed for the final report

% Summarise purpose of and motivation for the document (problem statement and scope) 

% Explain the final solution and link back to problem 

%Summarise the most important findings 

A useful experiment within cellular biology is long term cell imaging. Getting time-lapse footage of cell growth requires adequate incubation to ensure that the cell culture is able to grow uninhibited, and under the desired conditions. Stage-top incubators aim to solve this problem by providing an incubation chamber that is small enough to mount directly on top of existing microscopes that are present in most labs. The main issue with existing stage-top incubators is that they do not provide adequate containment capabilities for microbes that are classified as potentially lethal, such as those present in a bio-safety level three laboratory. 

This project aims to explore the possibility of developing a custom stage top incubator that is capable of functioning whilst providing suitable containment that allows these long term observations to occur within a bio-safety level three laboratory without the necessity for external containment.

This report documents the process of creating a conceptual design and prototype to demonstrate the bio-containment aspect of this project in order to determine the feasibility of taking on the full project for stage top incubation. For this purpose, a chamber is designed to house a Greiner 96 well microplate, and provide the capability to hermetically seal the multi-well plate inside, whilst still allowing for microscopy.

The design went through three major iterations, whereby feedback was used to alter and improve the design to a point whereby it could be 3D printed as a proof of concept.

Ultimately the conceptual design indicated that it would be possible to manufacture a bio-containment chamber that would be suitable for use within a BSL-3 lab. The design could be further improved upon to allow for environmental control. This shows that it would be feasible to undertake the project of designing a stage-top incubator with basic environmental control. 
% =====================================================
\section{Recommendations}
% Future recommendations for improvements in your project
The 3D printed model provided some insight into the performance of the bio-containment chamber. Whilst the prototype does successfully demonstrate the functionality of the chamber, it does indicate that there are areas where improvements could be made.

Whilst the stand-offs designed into the lid, that are used to press the multi-well plate into the gasket on the base, do provide enough force to prevent the plate from moving, the design does not take into consideration other multi-well plate manufactures, and so is isolated to working with the Greiner multi-well plate. Additionally it does not take into consideration the variance in plate heights between units. Further, with the primary mounting force originating in the four corners of the chamber, this allows for some elastic deformation at the stand-off positions which reduces the mounting force. 

% ----------------------------------------------------
\ifstandalone
\bibliography{../Bibliography/References.bib}
\fi
\end{document}
% ----------------------------------------------------